%% -*- coding: utf-8 -*-
\documentclass[12pt,a4paper]{scrartcl} 
\usepackage[utf8]{inputenc}
\usepackage[english,russian]{babel}
\usepackage{indentfirst}
\usepackage{misccorr}
\usepackage{graphicx}
\usepackage{amsmath}
\usepackage{listings}
\usepackage{xcolor}
\usepackage{float}
\lstset{
  basicstyle=\ttfamily\small,
  keywordstyle=\color{blue},
  commentstyle=\color{gray},
  stringstyle=\color{green},
  numberstyle=\tiny\color{gray},
  stepnumber=1,
  numbersep=10pt,
  backgroundcolor=\color{white},
  showspaces=false,
  showstringspaces=false,
  showtabs=false,
  frame=none, 
  tabsize=2,
  captionpos=b,
  breaklines=true,
  breakatwhitespace=false,
  title=\lstname
}
\begin{document}
	\begin{titlepage}
		\begin{center}
			\large
			МИНИСТЕРСТВО НАУКИ И ВЫСШЕГО ОБРАЗОВАНИЯ РОССИЙСКОЙ ФЕДЕРАЦИИ
			
			Федеральное государственное бюджетное образовательное учреждение высшего образования
			
			\textbf{АДЫГЕЙСКИЙ ГОСУДАРСТВЕННЫЙ УНИВЕРСИТЕТ}
			\vspace{0.25cm}
			
			Инженерно-физический факультет
			
			Кафедра автоматизированных систем обработки информации и управления
			\vfill
			
			\textsc{Отчет по практике}\\[5mm]
			
			{\LARGE Быстрая сортировка и сортировка методом слияния \textit{}}
			\bigskip
			
			3 курс, группа УТС
		\end{center}
		\vfill
		
		\newlength{\ML}
		\settowidth{\ML}{«\underline{\hspace{0.7cm}}» \underline{\hspace{2cm}}}
		\hfill\begin{minipage}{0.5\textwidth}
			Выполнил:\\
			\underline{\hspace{\ML}} Г.\,К.~Арустамов\\
			«\underline{\hspace{0.7cm}}» \underline{\hspace{2cm}} 2025 г.
		\end{minipage}%
		\bigskip
		
		\hfill\begin{minipage}{0.5\textwidth}
			Руководитель:\\
			\underline{\hspace{\ML}} С.\,В.~Теплоухов\\
			«\underline{\hspace{0.7cm}}» \underline{\hspace{2cm}} 2025 г.
		\end{minipage}%
		\vfill
		
		\begin{center}
			Майкоп, 2025 г.
		\end{center}
	\end{titlepage}

    \newpage

\section{Введение}
\label{sec:intro}

\subsection{Текстовая формулировка задачи (Вариант 4)}
Сортировки Быстрая и Слиянием.

\subsection{Теория метода}

"Быстрая сортировка", хоть и была разработана более 40 лет назад, является наиболее широко применяемым и одним их самых эффективных алгоритмов. Метод основан на подходе "разделяй-и-властвуй". Общая схема такова:

\begin{itemize}
  \item из массива выбирается некоторый опорный элемент a[i],
  \item запускается процедура разделения массива, которая перемещает все ключи, меньшие, либо равные a[i], влево от него, а все ключи, большие, либо равные a[i] - вправо,
  \item теперь массив состоит из двух подмножеств, причем левое меньше, либо равно правого,
  \item для обоих подмассивов: если в подмассиве более двух элементов, рекурсивно запускаем для него ту же процедуру.
\end{itemize}

В конце получится полностью отсортированная последовательность.


\section{Ход работы}


\begin{verbatim}
	
	"Быстрая сортировка"
	
	#include <iostream>
	
	namespace std {
		int partition(int arr[], int low, int high) {
			int pivot = arr[high]; 
			int i = (low - 1); 
			
			for (int j = low; j <= high - 1; j++) {
				if (arr[j] <= pivot) {
					i++;
					int temp = arr[i];
					arr[i] = arr[j];
					arr[j] = temp;
				}
			}
			int temp = arr[i + 1];
			arr[i + 1] = arr[high];
			arr[high] = temp;
			
			return (i + 1);
		}
		
		void quickSort(int arr[], int low, int high) {
			if (low < high) {
				int pi = partition(arr, low, high);
				
				quickSort(arr, low, pi - 1);
				quickSort(arr, pi + 1, high);
			}
		}
		
		void quickSort(int arr[], int size) {
			quickSort(arr, 0, size - 1);
		}
	}
	
	int main() {
		using namespace std;
		setlocale(0, "ru");
		int arr[] = { 10, 7, 8, 9, 1, 5 };
		int n = sizeof(arr) / sizeof(arr[0]);
		
		cout << "Исходный массив: ";
		for (int i = 0; i < n; i++)
		cout << arr[i] << " ";
		cout << "\n";
		
		quickSort(arr, n);
		
		cout << "Отсортированный массив: ";
		for (int i = 0; i < n; i++)
		cout << arr[i] << " ";
		cout << "\n";
		
		return 0;
	}
	
	
	"Сортировка методом слияния"
	
	#include <iostream>
	
	namespace std {
		void merge(int arr[], int left, int mid, int right) {
			int n1 = mid - left + 1;
			int n2 = right - mid;
			int* L = new int[n1];
			int* R = new int[n2];
			for (int i = 0; i < n1; i++)
			L[i] = arr[left + i];
			for (int j = 0; j < n2; j++)
			R[j] = arr[mid + 1 + j];
			int i = 0, j = 0, k = left;
			while (i < n1 && j < n2) {
				if (L[i] <= R[j]) {
					arr[k] = L[i];
					i++;
				}
				else {
					arr[k] = R[j];
					j++;
				}
				k++;
			}
			while (i < n1) {
				arr[k] = L[i];
				i++;
				k++;
			}
			while (j < n2) {
				arr[k] = R[j];
				j++;
				k++;
			}
			delete[] L;
			delete[] R;
		}
		void mergeSort(int arr[], int left, int right) {
			if (left < right) {
				int mid = left + (right - left) / 2;
				
				mergeSort(arr, left, mid);
				mergeSort(arr, mid + 1, right);
				
				merge(arr, left, mid, right);
			}
		}
		void mergeSort(int arr[], int size) {
			mergeSort(arr, 0, size - 1);
		}
	}
	
	int main() {
		using namespace std;
		setlocale(0, "ru");
		int arr[] = { 12, 11, 13, 5, 6, 7 };
		int n = sizeof(arr) / sizeof(arr[0]);
		
		cout << "Исходный массив: ";
		for (int i = 0; i < n; i++)
		cout << arr[i] << " ";
		cout << "\n";
		
		mergeSort(arr, n);
		
		cout << "Отсортированный массив: ";
		for (int i = 0; i < n; i++)
		cout << arr[i] << " ";
		cout << "\n";
		
		return 0;
	}
	
	
\end{verbatim}


\section{Скриншоты программы}
\label{sec:program-shots}

На следующих изображениях представлены примеры интерфейса реализованной программы:

\begin{figure}[H]
    \centering
    \includegraphics[width=0.7\linewidth]{Быстрая сортировка.png}
    \caption{Результат быстрой сортировки}
    \label{fig:desktop-view}
\end{figure}

\begin{figure}[H]
    \centering
    \includegraphics[width=0.7\linewidth]{Сортировка слиянием.png}
    \caption{Результат сортировки слиянием}
    \label{fig:mobile-view}
\end{figure}

\section{Доступ к исходному коду и результатам работы}
\text {Репозиторий:} 
\texttt{https://github.com/grigorijarustamov-alt/pract1.git}
\end{document}